\chapter{规范理论}\label{chap:1}
\begin{refsection}
本导论性章节旨在为后续章节的数学发展提供总体背景说明。我们给出全文所研究的基本方程，并简要介绍它们在物理学中的应用背景。

\section{Yang--Mills--Higgs理论}\label{sec:ymh}

在具有 $d$ 个空间维度的\iemph{Yang--Mills--Higgs}理论中，其动力学变量包括一个\iemph{规范势}（\iemph{联络}）
\[
A = A_j(x)\,dx^j
\]
以及一个\iemph{标量（Higgs）场}
\[
\Phi = \Phi(x) .
\]

在一个时间维加 $d$ 个空间维的情形下，变量形式类似。分量 $A_j(x)$ 取值于一个矩阵李群 $G$ 的李代数 $\mathfrak{g}$，该群 $G$ 被称为\iemph{规范群}。场 $\Phi$ 取值于一个向量空间 $\mathbb{L}$（称为\iemph{内部对称空间}），群 $G$ 作为变换群作用于其上。

Yang--Mills势定义了一个\iemph{场}（\iemph{曲率}）
\begin{equation}\label{eq:1.1.1}
F = dA + A\wedge A = \frac12 F_{ij}(x)\,dx^i\wedge dx^j ,
\end{equation}
其分量为
\begin{equation}\label{eq:1.1.2}
F_{ij}(x)
 = \partial_i A_j(x) - \partial_j A_i(x)
   + [A_i(x),A_j(x)] .
\end{equation}
此处及以下均采用求和约定。

Higgs场通过协变导数与联络发生耦合，这称为“最小耦合”。设 $\rho$ 为群 $G$ 在 $\mathbb{L}$ 上的一个线性表示，则 $\rho$ 诱导出李代数 $\mathfrak{g}$ 在 $\mathbb{L}$ 上的表示，\iemph{协变导数}定义为
\begin{subequations}
  \begin{equation}\label{eq:1.1.3a}
    (\nabla_{A})_j(\Phi) = \nabla_j\Phi + \rho(A_j)(\Phi),
  \end{equation}
  并且
  \begin{equation}\label{eq:1.1.3b}
    D_A\Phi = (\nabla_A)_j(\Phi)\,dx^j .
  \end{equation}
\end{subequations}
在\cref{chap:1}中，我们研究定义在 $\mathbb{R}^2$ 上的\uwave{阿贝尔Higgs模型}\index{Higgs模型!阿贝尔|hyperpage}，其规范群为 $G=U(1)$。
\begin{rmk}
文献引用示例：Belavin 等人的 BPST 瞬子 \cite{belavin1975}；Polyakov 的工作 \cite{polyakov1974,polyakov1977}。
\end{rmk}

\begin{thm}[示例定理]\label{thm:example}
设 $G$ 为紧 Lie 群，$A$ 为 $\mathfrak{g}$-值联络。则 $F_A = dA + A \wedge A$ 满足 Bianchi 恒等式 $D_A F_A = 0$。
\end{thm}

\begin{proof}[Bianchi 恒等式]
直接计算协变外微分即可。
\end{proof}

\printbibliography[heading=subbibliography,title=本章参考文献]
\end{refsection}
